\documentclass[11pt]{article}
\usepackage[a4paper,margin=1in]{geometry}
\usepackage{amsmath,amssymb,amsthm}
\usepackage{array}
\usepackage{hyperref}

\newcommand{\dd}{\mathrm{d}}
\newcommand{\bx}{\mathbf{x}}
\newcommand{\bv}{\mathbf{v}}
\newcommand{\bk}{\mathbf{k}}
\newcommand{\eps}{\varepsilon}
\newcommand{\diel}{\epsilon}
\newcommand{\avg}[1]{\langle #1 \rangle}

\theoremstyle{definition}
\newtheorem{target}{Target statement}
\newtheorem{obstruction}{Obstruction}

\title{Extending the Deng--Hani--Ma closure analysis to plasmas:\\
what carries over, what breaks, and what to compute}
\author{Working note}
\date{16 November 2026}

\begin{document}
\maketitle

\noindent\emph{Status and scope.} This note does not prove a theorem. A
long-time derivation of a collisional plasma kinetic equation from Coulomb
particle dynamics is an open problem, and nothing here closes it. The note
does three things: it maps each ingredient of the Deng--Hani--Ma analysis to
its plasma counterpart; it identifies where the proof strategy fails and what
would have to replace it; and it turns the result into a computational track
in the format of \texttt{math\_to\_algorithm.md}. The formal derivation in
Section~\ref{sec:formal} is the classical Bogoliubov argument, not a rigorous
one. The literature summary in Section~\ref{sec:known} was checked on the
date above.

\tableofcontents

\section{What is being extended}

Stripped of technique, the hard-sphere analysis has six components.
\begin{enumerate}
  \item A reversible microscopic system: $N$ hard spheres of diameter $\eps$.
  \item A scaling limit with one small parameter: Boltzmann--Grad,
    $N\eps^{d-1} = \alpha$ fixed, $\eps\to 0$.
  \item A retained state, the one-particle density $f_1$, and a closure,
    $f_2 \approx f_1\otimes f_1$ for particles about to collide.
  \item A cumulant expansion of $f_s$ in which the connected parts are the
    closure error, each represented as a sum over collision histories
    (molecules).
  \item A gain: histories with cycles (recollisions) are suppressed by a
    positive power of $\eps$, and a cutting algorithm proves this uniformly.
  \item Time layering: the cumulant ansatz is propagated across many short
    layers, reaching any time on which the Boltzmann solution exists.
\end{enumerate}
The plasma extension must supply an analogue of each. Components 1--3 have
clean analogues. Components 4--6 do not, and the reasons are structural.

\section{The plasma reduction ladders}

Take $N$ particles of unit mass in a periodic box, interacting through a pair
potential $\phi$ in the mean-field scaling,
\begin{equation}
  \dot{\bx}_i = \bv_i, \qquad
  \dot{\bv}_i = -\frac{1}{N}\sum_{j\ne i}\nabla\phi(\bx_i - \bx_j).
  \label{eq:nbody}
\end{equation}
For a one-component plasma $\phi$ is the Coulomb potential in a neutralizing
background; a Yukawa potential $\phi(r) = e^{-\kappa r}/r$ models a screened
or dusty plasma. In physical units the expansion parameter is the inverse
plasma parameter $1/\Lambda$, $\Lambda\sim n\lambda_D^3$, which plays the
role of $1/N$ here.

\begin{center}
\renewcommand{\arraystretch}{1.45}
\begin{tabular}{>{\raggedright}p{0.6cm}>{\raggedright}p{3.6cm}>{\raggedright}p{3.4cm}>{\raggedright}p{2.6cm}p{3.6cm}}
  & Transition & Limit & Time scale & Hard-sphere analogue \\ \hline
  P1 & particles $\to$ Vlasov & $N\to\infty$ & $O(1)$
     & none: no collisions at leading order \\
  P2 & particles $\to$ Landau or Lenard--Balescu & $N\to\infty$
     & $O(N)$ & hard spheres $\to$ Boltzmann \\
  P3 & Vlasov--Landau $\to$ fluid & $\mathrm{Kn}\to 0$ & fluid
     & Boltzmann $\to$ Euler, Navier--Stokes--Fourier \\
  P4 & Vlasov $\to$ quasilinear, weak turbulence & small amplitude, random phases
     & $O(\text{amplitude}^{-2})$ & NLS $\to$ wave kinetic equation \\
\end{tabular}
\end{center}

The essential difference is already visible in P1 and P2. For hard spheres
there is one kinetic limit. For a plasma there are two nested ones: the
leading-order limit is collisionless, and the collisional equation appears
only as a $1/N$ correction that matters on times of order $N$.

\section{Formal derivation of the collisional closure}
\label{sec:formal}

\subsection{Hierarchy and cluster expansion}

The first BBGKY equation for \eqref{eq:nbody} is
\begin{equation}
  \partial_t f_1 + \bv\cdot\nabla_{\bx} f_1
  = \frac{N-1}{N}\int \nabla\phi(\bx - \bx')\cdot\nabla_{\bv}
    f_2(z, z')\,\dd z', \qquad z = (\bx,\bv).
\end{equation}
Write the cluster expansion with the expected orders made explicit,
\begin{equation}
  f_2 = f\otimes f + \frac{1}{N}\,g, \qquad
  f_3 = f^{\otimes 3} + \frac{1}{N}\sum f\otimes g + \frac{1}{N^2}\,h ,
\end{equation}
with $f = f_1$. Then
\begin{equation}
  \partial_t f + \bv\cdot\nabla_{\bx} f
  - (\nabla\phi * \rho_f)\cdot\nabla_{\bv} f
  = \frac{1}{N}\,\nabla_{\bv}\cdot\!\int \nabla\phi(\bx-\bx')\,
    g(z,z')\,\dd z' + O\!\left(\frac{1}{N}\right)_{\text{self}} .
  \label{eq:f1}
\end{equation}
The left side is the Vlasov operator. The right side is the collision term,
and it is carried entirely by the two-point cumulant $g$.

At leading order $g$ obeys a linear equation driven by the product state,
\begin{equation}
  \partial_t g + (\mathcal{L}_1 + \mathcal{L}_2)\,g
  = \nabla\phi(\bx_1 - \bx_2)\cdot(\nabla_{\bv_1} - \nabla_{\bv_2})
    \,f(z_1) f(z_2) + O\!\left(\frac{1}{N}\right),
  \label{eq:g}
\end{equation}
where $\mathcal{L}_i$ is the Vlasov operator linearized about $f$, acting on
variable $i$, and the $O(1/N)$ remainder contains the three-point cumulant
$h$.

\subsection{Bogoliubov closure}

Two assumptions close the system.
\begin{itemize}
  \item \textbf{Truncation:} drop $h$.
  \item \textbf{Time-scale separation:} $g$ relaxes on the plasma period,
    $f$ evolves on the collision time, which is $N$ times longer. So solve
    \eqref{eq:g} with $f$ frozen and take $t\to\infty$; $g$ becomes a
    functional of the current $f$, with no memory of its own initial data.
\end{itemize}
For spatially homogeneous $f(\bv)$ the result is the Lenard--Balescu
equation,
\begin{equation}
  \partial_t f = \frac{1}{N}\,\nabla_{\bv}\cdot\!\int
    \mathsf{B}(\bv,\bv_*;f)\,(\nabla_{\bv} - \nabla_{\bv_*})
    \,f(\bv) f(\bv_*)\,\dd\bv_* ,
  \label{eq:LB}
\end{equation}
\begin{equation}
  \mathsf{B}(\bv,\bv_*;f) = c_d\int
    \frac{\bk\otimes\bk\;|\hat\phi(\bk)|^2}
         {|\diel(\bk,\bk\cdot\bv)|^2}\;
    \delta\bigl(\bk\cdot(\bv - \bv_*)\bigr)\,\dd\bk ,
  \qquad
  \diel(\bk,\omega) = 1 + \hat\phi(\bk)\int
    \frac{\bk\cdot\nabla_{\bv} f}{\omega + i0 - \bk\cdot\bv}\,\dd\bv ,
\end{equation}
with $c_d$ a constant fixed by the Fourier convention. Setting $\diel = 1$
gives the Landau equation.

Three features matter for what follows.
\begin{enumerate}
  \item The delta function is a wave--particle resonance: two particles
    exchange momentum through the Fourier mode that is stationary in both
    their frames.
  \item The dielectric function $\diel$ contains the same resonant
    denominator, Landau contour and Plemelj structure as the plasma
    dispersion function in the Hammett--Perkins note. Screening is the
    linearized Vlasov response of all other particles to the interacting
    pair.
  \item For the bare Coulomb potential, $|\hat\phi|^2\propto k^{-4}$ and the
    $\bk$ integral diverges logarithmically at large $k$ if $\diel=1$ and is
    cut off by $\diel$ at small $k$. The large-$k$ divergence signals close
    encounters, where the weak-deflection expansion is invalid and a
    Boltzmann-type treatment is required. The Coulomb logarithm is the
    result of matching these regimes.
\end{enumerate}

\section{Dictionary}

\begin{center}
\renewcommand{\arraystretch}{1.5}
\begin{tabular}{>{\raggedright}p{3.1cm}>{\raggedright}p{5.2cm}p{6.2cm}}
  Ingredient & Hard spheres (Deng--Hani--Ma) & Plasma \\ \hline
  Small parameter & $\eps$, with $N\eps^{d-1}$ fixed
    & $1/N$ (physically $1/\Lambda$) \\
  Leading-order closure & $f_2 = f\otimes f$ gives the collision integral
    & $f_2 = f\otimes f$ gives the mean field; no collisions \\
  Role of the two-point cumulant & closure \emph{error}
    & the collision operator itself, eq.~\eqref{eq:f1} \\
  Closure error & all connected cumulants
    & three-point cumulant $h$, memory of $g$, close encounters \\
  Elementary event & instantaneous binary collision, $O(1)$ deflection, rare
    & every pair interacts at all times, $O(1/N)$ deflection each \\
  History object & collision graph (molecule), sparse
    & no discrete events; complete weighted graph, or perturbative diagrams \\
  Propagator between events & free transport
    & linearized Vlasov flow about $f$ (dressed, includes $\diel$) \\
  Cycles & recollisions, each costs a power of $\eps$
    & ring diagrams, all of the same order; resummed into $1/|\diel|^2$ \\
  Source of irreversibility & pre-collisional chaos
    & loss of memory in $g$ (Bogoliubov) after reversible phase mixing \\
  Stability of the limit equation & H-theorem
    & Vlasov has no H-theorem; needs linear (Penrose) stability of $f$ \\
\end{tabular}
\end{center}

\section{Where the proof strategy breaks}

\begin{obstruction}[No sparse collision graph]
The molecule expansion relies on collisions being discrete and rare, so that
a typical history is a tree and each cycle can be charged a power of $\eps$.
In \eqref{eq:nbody} every pair interacts continuously. The natural
replacement is a Duhamel expansion in the pair interaction, with diagrams in
place of molecules. The counting problem is then different in kind: the
number of diagrams at a given order is not controlled by rarity of events.
\end{obstruction}

\begin{obstruction}[Cycles are not small]
For hard spheres, cycles are the error. For a plasma an infinite class of
cyclic diagrams (rings) contributes at leading order to the collision
operator and must be summed exactly; the sum is the factor $1/|\diel|^2$. A
cutting argument that discards cycles would discard screening. Any extension
needs a resummation step before the combinatorial estimate, which amounts to
expanding around the linearized Vlasov flow instead of free transport.
\end{obstruction}

\begin{obstruction}[The background must be stable]
The dressed propagator is the linearized Vlasov semigroup about the current
$f$. It decays (Landau damping) only if $f$ is linearly stable; near an
unstable profile it grows exponentially and no cumulant bound can propagate.
Hard spheres have no counterpart to this. A theorem can be expected only for
$f$ that remain in a stable class, for example near a homogeneous
Penrose-stable equilibrium.
\end{obstruction}

\begin{obstruction}[Memory and echoes]
Landau damping is reversible phase mixing: information moves to fine
velocity scales and can return (plasma echo). The Bogoliubov step asserts
that $g$ forgets its past on the fast time scale. Over the long kinetic time
this must hold uniformly, which is the plasma form of the long-time problem
that layering solves for hard spheres. The layered cumulant ansatz would
need to carry fine-scale velocity regularity, not only size.
\end{obstruction}

\begin{obstruction}[Three collision regimes]
With Coulomb interaction the kinetic operator is assembled from close
encounters (Boltzmann-like), weak binary deflections (Landau) and collective
response (Lenard--Balescu). Only the first resembles the hard-sphere setting,
and it is subdominant by a factor $1/\ln\Lambda$. A single expansion does not
cover all three.
\end{obstruction}

\begin{obstruction}[Singularity of the force]
Even the leading-order limit P1 requires control of the Coulomb singularity
at short range, a difficulty with no hard-sphere analogue.
\end{obstruction}

\section{What a plasma theorem would say}

These are targets, stated to fix what the extension aims at. None is claimed.

\begin{target}[Cumulant hierarchy on the kinetic time scale]
Let $\phi$ be smooth and $f$ stay in a neighbourhood of a homogeneous
Penrose-stable equilibrium. Then for $\tau = t/N$ in a fixed interval
$[0,\tau_*]$ the connected cumulants of the $N$-particle system satisfy, in a
suitable weak norm, $\|g_s\| \lesssim N^{-(s-1)}$ uniformly in $\tau$.
\end{target}

\begin{target}[Lenard--Balescu limit]
Under the same hypotheses, the velocity marginal of $f_1$ at time $t = N\tau$
converges as $N\to\infty$ to the solution of \eqref{eq:LB} with time
$\tau$, for as long as that solution exists and remains in the stable class.
\end{target}

\begin{target}[Interpolation to the Boltzmann--Grad regime]
For the Yukawa potential with screening length $1/\kappa$, there is a family
of scalings interpolating between mean-field ($\kappa$ fixed, weak coupling)
and Boltzmann--Grad (short range, strong deflection, low density), with
Landau, Lenard--Balescu and Boltzmann operators as limits. The hard-sphere
theorem is the short-range endpoint.
\end{target}

The third target is not new: Nota, Vel\'azquez and Winter formulate the
interpolation as conjectures, with Boltzmann, Landau and Lenard--Balescu
limits selected by the scaling of the potential \cite{NVW}. It is the most
useful for this project: it places the
Deng--Hani--Ma result and the plasma on one axis, and dusty plasmas sit on
that axis at strong coupling, where none of the limits applies.

\section{What is known}
\label{sec:known}

This section was checked against the literature on 16 November 2026.
Citations were confirmed online unless marked otherwise in the bibliography;
several are preprints and are identified as such.

\paragraph{The reference point.} Deng, Hani and Ma derive the Boltzmann
equation from hard spheres in $\mathbb{R}^d$, $d\ge 2$, in the Boltzmann--Grad
limit, for as long as a suitable strong solution of the Boltzmann equation
exists \cite{DHM1}; the fluid limits are obtained on the torus in $d=2,3$ by
taking the collision rate to infinity after the Boltzmann--Grad limit
\cite{DHM2}. The short-time theory is due to Lanford \cite{Lanford} and
Gallagher, Saint-Raymond and Texier \cite{GST}. The wave-kinetic analogue is
\cite{DH-WKE,DH-long}.

\paragraph{P1, particles to Vlasov.}
\begin{itemize}
  \item \emph{Lipschitz forces:} classical \cite{BraunHepp,Dobrushin}, with a
    stability constant $e^{Ct}$, so the estimate is useful only for
    $t \ll \log N$.
  \item \emph{Longer times near stable states:} in one dimension with
    smooth potentials, mean-field validity up to times $N^{1/8-}$ near
    stable homogeneous profiles \cite{CagliotiRousset}.
  \item \emph{Singular forces weaker than Coulomb:} $|F|\sim |x|^{-\alpha}$
    with $\alpha < 1$, without cutoff \cite{HaurayJabin}.
  \item \emph{Coulomb with an $N$-dependent cutoff:} in three dimensions
    \cite{LazaroviciPickl}.
  \item \emph{Modulated energy:} covers first-order Coulomb flows, and
    second-order dynamics only for monokinetic data \cite{Serfaty}.
  \item \emph{Coulomb without cutoff:} a 2026 preprint claims the
    two-dimensional case for chaotic data and regular mean-field solutions
    \cite{DuerinckxJabin}. The three-dimensional case is open.
\end{itemize}

\paragraph{P2, particles to Landau or Lenard--Balescu.} Open. The available
results are:
\begin{itemize}
  \item \emph{Consistency for Landau:} the leading-order truncation of the
    weak-coupling expansion agrees with the Landau collision term; no
    convergence for positive times \cite{BPS}.
  \item \emph{From the truncated hierarchy:} convergence to Landau when the
    truncation of BBGKY is taken as a hypothesis \cite{VelazquezWinter,Winter}.
  \item \emph{Lenard--Balescu correction:} for a tagged particle in an
    equilibrium background with smooth interaction on the torus, the $1/N$
    correction is the linearized Lenard--Balescu operator, on times $N^r$ for
    a small $r > 0$, far short of the relaxation time $t\sim N$
    \cite{DuerinckxSR}. The same paper remarks that in one dimension the
    operator vanishes and the relevant scale is $N^2$.
  \item \emph{Cumulant bounds:} sharp estimates on correlation functions on
    the mean-field time scale for smooth interactions \cite{DuerinckxChaos}.
    These are the plasma counterpart of the cumulant bounds in component 4,
    but only on short times.
  \item \emph{Toy model:} convergence on $t\sim N$ to a linearized Landau
    equation for a truncated hierarchical model, not the particle system;
    its introduction states that a rigorous justification for the true
    dynamics ``is still beyond reach'' \cite{DuerinckxLeBihan}.
\end{itemize}
The first two target statements of the previous section therefore lie beyond
the current state of the art even in the linearized equilibrium setting.

\paragraph{Stability input.} Nonlinear Landau damping holds near linearly
stable homogeneous equilibria for analytic and Gevrey data
\cite{MouhotVillani,BMM}. In finite (Sobolev) regularity, plasma echoes
obstruct results of this type \cite{BedrossianEcho}. This is the rigorous
form of the fourth obstruction.

\paragraph{P3, kinetic to fluid.} A global Hilbert expansion gives
Euler--Poisson from Vlasov--Poisson--Boltzmann for small irrotational data
\cite{GuoJang}; conditional convergence to incompressible viscous
electro-magneto-hydrodynamics is in \cite{ArsenioSR}. These start from a
kinetic equation, not from particles, so no plasma analogue of the chained
particle-to-fluid result \cite{DHM2} exists.

\paragraph{P4, quasilinear and weak turbulence.} A velocity-diffusion limit
is proved for a prescribed stochastic field, and the authors state that a
complete proof of quasilinear theory for the self-consistent problem is out
of reach \cite{BardosBesse}. There is no plasma analogue of
\cite{DH-WKE,DH-long}.

\paragraph{One dimension.} Numerical relaxation on a time scale $\propto N^2$
was reported by Dawson \cite{Dawson}; kinetic equations at order $1/N^2$ have
been derived formally \cite{Fouvry}. The vanishing concerns the
self-consistent evolution of the whole distribution; a test particle in a
fixed bath still diffuses at order $1/N$.

\section{A fluctuation formulation}
\label{sec:fluct}

The remaining sections go beyond the dictionary and propose a strategy. The
scaling arguments are heuristic and are labelled as such. Much of the
picture turns out to be in the recent literature; Section~\ref{sec:prior}
records what is prior art and what is proposed here.

For smooth $\phi$ the empirical measure
$\mu_N = N^{-1}\sum_i \delta_{z_i(t)}$ is an exact weak solution of the Vlasov
equation. Put $f = \mathbb{E}\mu_N$ (which is exactly $f_1$) and
$\eta = \sqrt{N}\,(\mu_N - f)$, and write $F[h] = -\nabla\phi * \rho_h$ for
the force generated by $h$. Then, with no truncation,
\begin{align}
  \partial_t f + \bv\cdot\nabla_{\bx} f + F[f]\cdot\nabla_{\bv} f
    &= -\frac{1}{N}\,\nabla_{\bv}\cdot\mathbb{E}\bigl[F[\eta]\,\eta\bigr],
    \label{eq:mean}\\
  \partial_t \eta + \mathcal{L}_f\,\eta
    &= -\frac{1}{\sqrt{N}}\,\nabla_{\bv}\cdot
       \Bigl(F[\eta]\,\eta - \mathbb{E}\bigl[F[\eta]\,\eta\bigr]\Bigr),
    \label{eq:fluct}
\end{align}
where $\mathcal{L}_f\eta = \bv\cdot\nabla_{\bx}\eta + F[f]\cdot\nabla_{\bv}\eta
+ F[\eta]\cdot\nabla_{\bv} f$ is the linearized Vlasov operator. The
covariance of $\eta$ is
$f\,\delta_{z=z'} - f\otimes f + (1 - 1/N)\,g$: shot noise plus the two-point
cumulant.

Equation \eqref{eq:fluct} is a wave equation with a quadratic nonlinearity of
strength $\alpha = N^{-1/2}$ and random initial data, and \eqref{eq:mean} says
the retained state evolves at rate $\alpha^2$. This is the scaling of the
wave-kinetic problem, where the kinetic time is $\alpha^{-2}$:
\begin{center}
\renewcommand{\arraystretch}{1.4}
\begin{tabular}{>{\raggedright}p{4.0cm}>{\raggedright}p{4.6cm}p{5.6cm}}
  & NLS $\to$ wave kinetic & Fluctuating Vlasov $\to$ Lenard--Balescu \\ \hline
  Field & $\hat u_k$ & $\eta_{\bk}(\bv)$ \\
  Nonlinearity strength & $\alpha$ & $N^{-1/2}$ \\
  Kinetic time & $\alpha^{-2}$ & $N$ \\
  Initial data & independent Gaussian modes & Poisson-type shot noise \\
  Mode frequency & $|k|^2$ & $\bk\cdot\bv$ (continuous in $\bv$) \\
  Resonance & four-wave manifold & $\bk\cdot(\bv - \bv_*) = 0$ \\
  Linear propagator & diagonal, unitary & linearized Vlasov: non-normal, damped \\
  Retained state & spectrum $n_k$ & $f$, which also sets the propagator \\
\end{tabular}
\end{center}
So ladder P2 is ladder P4 with the wave spectrum supplied by particle
discreteness, and the natural template for a proof is the wave-kinetic
analysis of Deng and Hani \cite{DH-WKE,DH-long}, not the hard-sphere one.
Four features have no counterpart there and are where new work is needed:
the nonlinearity loses a velocity derivative; the propagator is not unitary;
the background $f$ evolves and feeds back; and the data are not Gaussian, with
cumulants of order $s$ of size $N^{-(s-2)/2}$, the same order as the
nonlinear corrections.

\section{Three time scales and two thresholds (heuristic)}
\label{sec:scales}

Work in the frame moving with free transport, where a fluctuation created at
time $s$ with wavenumber $\bk$ has velocity-frequency $\bk\,(t-s)$ at time
$t$. A velocity derivative acting on it costs a factor $|\bk|(t-s)$. The
field $F[\eta]$ is a stationary random field of size $O(1)$ with correlation
time $O(1)$.

\begin{itemize}
  \item \textbf{Worst-case bound.} One Duhamel iterate of \eqref{eq:fluct}
    has relative size at most
    $N^{-1/2}\int_0^t |F|\,k s\,\dd s \sim N^{-1/2}t^2$. A deterministic
    estimate therefore fails at $t\sim N^{1/4}$.
  \item \textbf{Random-phase bound.} Because $F$ decorrelates in $O(1)$
    time, the same iterate has variance
    $N^{-1}\int_0^t k^2 s^2\,\dd s \sim t^3/N$. An estimate that uses this
    cancellation fails at $t\sim N^{1/3}$.
  \item \textbf{What happens at $N^{1/3}$.} The breakdown is physical. A
    particle's velocity diffuses with coefficient $D\sim 1/N$, so its
    position departs from the unperturbed orbit by
    $\avg{|\Delta x|^2} = \tfrac{2}{3}Dt^3$. A phase-mixed fluctuation obeys,
    in the moving frame, $\dot a = -D k^2 t^2 a$, hence
    \begin{equation}
      a(t) = a(0)\,\exp\!\left(-\tfrac{1}{3}\,D k^2 t^3\right),
      \qquad
      \tau_{\mathrm{mem}} \sim \left(\frac{N}{k^2}\right)^{1/3}.
      \label{eq:mem}
    \end{equation}
\end{itemize}

\begin{center}
\renewcommand{\arraystretch}{1.4}
\begin{tabular}{>{\raggedright}p{2.6cm}>{\raggedright}p{5.2cm}p{6.4cm}}
  Time & Meaning & Hard-sphere counterpart \\ \hline
  $O(1)$ & correlation time of the field; $g$ relaxes to its Bogoliubov form
     & collision duration (zero) \\
  $N^{1/3}$ & memory time: orbits decorrelate, phase-mixed information is erased
     & none \\
  $N$ & relaxation of $f$ & mean free time \\
\end{tabular}
\end{center}

None of these scales is new. The cubic exponential \eqref{eq:mem} is
Dupree's resonance broadening \cite{Dupree} and the collisional attenuation
of echoes \cite{ONeil}. In the derivation problem, Duerinckx and Le Bihan
state that naive resolvent bounds control their expansion ``at best'' for
$t\ll N^{1/3}$ and that renormalization turns free propagators into
hypoelliptic ones, which compensates for filamentation
\cite{DuerinckxLeBihan}; Bodineau and Le Bris find a Gr\"onwall barrier at
$N^{1/4}$, extended towards $N^{1/3}$ by averaging of fluctuations, for a
tagged particle in a free gas \cite{BodineauLeBris}. What this section adds
is the reading of the two exponents as worst-case versus random-phase bounds
on the single term $N^{-1/2}\nabla_{\bv}\cdot(F[\eta]\eta)$.

The Markovian closure is correct because $N^{1/3}\ll N$: memory is erased
long before $f$ changes. The difficulty of the proof is that the erasure is
itself a collisional effect, of the same origin as the equation being
derived.

\section{Proposed techniques}
\label{sec:techniques}

\subsection{T1. Expand the fluctuation, not the hierarchy}

Apply a Duhamel expansion to \eqref{eq:fluct} about the linearized Vlasov
flow and organize the terms as in the wave-kinetic theory: trees for the
iterates of $\eta$, paired into couples when the expectation in
\eqref{eq:mean} is taken.
\begin{itemize}
  \item The dressed propagator $e^{-t\mathcal{L}_f}$ is explicit for
    homogeneous $f$: free transport plus a rank-one correction in $\bv$
    governed by $1/\diel$. Using it removes the ring diagrams from the
    combinatorics altogether, which disposes of the second obstruction.
  \item The leading couples reproduce \eqref{eq:LB}; the resonance
    $\delta(\bk\cdot(\bv-\bv_*))$ arises from the time integral of
    $e^{-i\bk\cdot(\bv-\bv_*)s}$ exactly as the resonant manifold does for
    waves.
  \item Pairing must use a calculus for independent, non-Gaussian particles
    in place of Wick's theorem. Duerinckx's Glauber calculus
    \cite{DuerinckxChaos} is the available tool; at present it is controlled
    only for $t\ll\log N$.
\end{itemize}
Fluctuation formulations of the Lenard--Balescu problem exist
\cite{Lancellotti,FeliachiBouchet}, and a Dyson expansion in Feynman diagrams
with renormalization has been carried out on the correlation hierarchy of a
toy model \cite{DuerinckxLeBihan}. A limited search found no application of
the random-data wave-kinetic method to \eqref{eq:fluct}; that is a null
search result, not a proof of absence. The expected gain is the one noted in
\cite{DuerinckxSR}: on the hierarchy each operator carries a velocity
derivative and summability seems to need analyticity, whereas mean-square
estimates over random data do not require regularity of the data.

\subsection{T2. Layer inside the memory window}

Re-expand at times $t_l = l\tau$ with
\begin{equation}
  1 \ll \tau \ll N^{1/3}.
\end{equation}
The lower bound lets $g$ reach its Bogoliubov form within a layer; the upper
bound keeps the expansion parameter $\tau^3/N$ small, so no resummation of
orbit perturbations is needed inside a layer. Term by term, contributions of
order $2m$ can be as large as $(\tau/N)(\tau^3/N)^{m-1}$; cancellations among
them are expected, as for irregular chains in the wave problem, but the
window already suffices for convergence.

The layered ansatz, playing the role of the cumulant ansatz of
Deng--Hani--Ma, would be
\begin{equation}
  g(t_l) = \mathcal{G}[f(t_l)] + r_l ,
\end{equation}
where $\mathcal{G}[f]$ is the Bogoliubov (dressed-particle) correlation and
$r_l$ is a remainder measured in a weak norm in the moving frame, with a
weight that records the age of each component, so that components older than
$\tau_{\mathrm{mem}}$ are discounted by the factor in \eqref{eq:mem}. Higher
cumulants are carried with bounds $N^{-(s-1)}$ and at most polynomial growth
in the moving-frame derivative.

\subsection{T3. Self-generated suppression of echoes}

The fourth obstruction is that phase-mixed information can return. The
proposal is that it cannot return after $\tau_{\mathrm{mem}}$, by a bootstrap.
\begin{itemize}
  \item \textbf{Mechanism.} Weak collisions of strength $\nu$ damp fine
    velocity structure at the enhanced rate $\nu^{1/3}$ and suppress echoes
    for small perturbations. For Vlasov--Poisson with Fokker--Planck or
    Landau collisions this is proved for perturbations of size
    $\ll\nu^{1/3}$ in high Sobolev regularity
    \cite{Bedrossian17,CLN}.
  \item \textbf{The problem is critical.} Here $\nu\sim 1/N$ and the
    fluctuation amplitude is $N^{-1/2}$. That is below the high-regularity
    threshold $\nu^{1/3} = N^{-1/3}$, but shot noise has no regularity, and
    in low regularity the sharp threshold is reported to be $\nu^{1/2}$
    \cite{ZhaoZi}, which is exactly $N^{-1/2}$. So discreteness noise sits
    at the borderline of deterministic echo suppression, with no room to
    spare. A deterministic argument cannot be expected to close; the margin
    has to come from independence of the particles.
  \item \textbf{Bootstrap.} Assume the layered ansatz up to $t_l$. It
    implies that a tagged particle diffuses in velocity at rate at least
    $c/N$ over the next layer. That in turn gives the damping
    \eqref{eq:mem} for the remainder, which closes the ansatz at $t_{l+1}$.
\end{itemize}
Two things make this harder than the collisional-Vlasov theorems. The
collisions are not an external operator but are produced by the fluctuations
being estimated. And, by the criticality just noted, smoothness must be
replaced by independence: the estimates have to be in mean square over the
initial data, as in random-data wave theory.

The toy-model theorem of \cite{DuerinckxLeBihan} reaches $t\sim N$ under an
added velocity diffusion of strength $\kappa/N$, which the authors say they
``cannot yet dispense with''. The bootstrap proposed here is a route to
removing that hypothesis: the diffusion would be supplied by the dynamics.

\subsection{T4. An intermediate theorem at $N^{1/3}$}

The bootstrap needs only a statement about a tagged particle on times just
past $N^{1/3}$, far short of the kinetic time.

\begin{target}[Orbit decorrelation]
For $\phi$ smooth, $f$ near a homogeneous Penrose-stable equilibrium, and
$1\ll t\le N^{1/3+\delta}$, a tagged particle satisfies
\[
  \mathbb{E}\exp\bigl(i\bk\cdot(\bx(t) - \bx(0) - \bv(0)t)\bigr)
  = \exp\!\left(-\frac{t^3}{3N}\,\bk\cdot\mathsf{D}(\bv(0))\,\bk\right)
    + o(1),
\]
with $\mathsf{D}$ the Lenard--Balescu diffusion matrix.
\end{target}

For a tagged particle in a non-interacting background in dimension
$d\ge 4$, convergence to the linear Landau diffusion on $t\sim N$ is proved
in \cite{BodineauLeBris}, which contains this statement in that setting. The
target is the interacting, screened case. It is the natural next step after
the tagged-particle result of \cite{DuerinckxSR}: it asks for the collisional correction on a time scale
where it is first visible, and it is the single input from which the memory
cutoff follows. It is also directly measurable (Section~\ref{sec:tests}).

\subsection{T5. Split the interaction at the interparticle spacing}

For Coulomb interaction in three dimensions, with lengths in units of the
Debye length, the $90^\circ$ deflection radius is $b_{90}\sim N^{-1}$ and the
interparticle spacing is $a\sim N^{-1/3}$. Write
$\phi = \phi_< + \phi_>$ with the split at a scale just below $a$.
\begin{itemize}
  \item Encounters through $\phi_<$ are binary and do not overlap in space
    and time. They form a sparse event graph, to which the molecule
    technique applies in a grazing-collision regime.
  \item $\phi_>$ acts on many particles at once and is treated by T1--T3.
\end{itemize}
Each range of impact parameters contributes to velocity diffusion in
proportion to its logarithmic width, so to leading logarithmic order
\begin{equation}
  \frac{\ln(a/b_{90})}{\ln(1/b_{90})} = \frac{2}{3}
  \quad\text{(binary)},
  \qquad
  \frac{\ln(1/a)}{\ln(1/b_{90})} = \frac{1}{3}
  \quad\text{(many-body)} .
\end{equation}
The split at the interparticle spacing is standard; I did not find the
fractions stated in the literature, and one review is reported to attribute
most of the transport to impact parameters between the spacing and the Debye
length, which would disagree. The fractions should be treated as a
leading-logarithm estimate to be checked, numerically and against that
review. On this accounting two thirds of the Coulomb logarithm is within
reach of a hard-sphere-type analysis, and only one third requires the fluctuation
theory. A hybrid expansion would use event histories for $\phi_<$ with the
dressed propagator of $\phi_>$ between events.

\subsection{Prior art and what is claimed}
\label{sec:prior}

Checked by a literature search on 16 November 2026. A null search result is
recorded as such and is not a claim of novelty.
\begin{center}
\renewcommand{\arraystretch}{1.4}
\begin{tabular}{>{\raggedright}p{5.2cm}p{9.0cm}}
  Ingredient & Status \\ \hline
  Exact fluctuation system & classical (Klimontovich); rigorous fluctuation
     theory in \cite{BraunHepp,Lancellotti} \\
  Analogy $\alpha = N^{-1/2}$, $T_{\mathrm{kin}} = \alpha^{-2}$; use of the
     random-data wave-kinetic method & not found in a limited search \\
  Diagram expansion with renormalization & done on the hierarchy for a toy
     model \cite{DuerinckxLeBihan} \\
  $N^{1/3}$ barrier; diffusion at rate $1/N$ cures filamentation & in
     \cite{DuerinckxLeBihan}; exponents $N^{1/4}$, $N^{1/3}$ also in
     \cite{BodineauLeBris} \\
  Link to enhanced dissipation and echo thresholds with $\nu = 1/N$;
     criticality of shot noise & not found stated; follows by combining
     \cite{Bedrossian17,CLN,ZhaoZi} with $\nu = 1/N$ \\
  Layering in the window $1\ll\tau\ll N^{1/3}$ with an age-weighted
     remainder & proposed here; layering itself is from
     \cite{DHM1,DH-long} \\
  Orbit-decorrelation target & proved for a free-gas background, $d\ge 4$
     \cite{BodineauLeBris}; open with interaction \\
  Interpolation between scalings & conjectured in \cite{NVW} \\
  Two-thirds, one-third split & elementary; not found stated; possibly in
     tension with a review, unchecked \\
\end{tabular}
\end{center}

\subsection{Proof architecture}

\begin{center}
\renewcommand{\arraystretch}{1.4}
\begin{tabular}{>{\raggedright}p{0.5cm}>{\raggedright}p{6.6cm}p{6.8cm}}
  & Step & Status \\ \hline
  1 & Exact fluctuation system \eqref{eq:mean}--\eqref{eq:fluct} & proved (\texttt{orbit\_decorrelation\_theorem.tex}, Lemma 1) \\
  2 & Dressed propagator and its decay for stable $f$ & known \\
  3 & Cumulant bounds on $O(1)$ times & known \cite{DuerinckxChaos} \\
  4 & One layer: expansion converges for $\tau\ll N^{1/3}$ and gives
      $(\tau/N)\,Q_{\mathrm{LB}}$ & proposed (T1, T2); needs mean-square
      estimates \\
  5 & Orbit decorrelation at $N^{1/3}$ & proved for a prescribed Markovian Gaussian field (\texttt{orbit\_decorrelation\_theorem.tex}); open for the particle system \\
  6 & Remainder cannot return: bootstrap & proposed (T3); open, the crux \\
  7 & Iterate over $N/\tau$ layers with $f$ in the stable class & needs
      stability of \eqref{eq:LB} in that class \\
  8 & Coulomb: hybrid expansion & proposed (T5); open \\
\end{tabular}
\end{center}

\section{Computational translation: Algorithm D}

This follows the layered closure state $S_l = (y_l, c_l, h_l, d_l,\sigma_l)$
of \texttt{math\_to\_algorithm.md}. The baseline reduced equation stays
explicit; a history-aware model predicts only the defect or a validity
indicator.

\subsection{Generator}

$N$-body dynamics \eqref{eq:nbody} in a periodic box with a Yukawa or
softened Coulomb potential, in two or three dimensions. Sweep $N$ at fixed
mean-field strength, and sweep the screening parameter $\kappa$.

One dimension needs care: for a homogeneous one-species system the
resonance condition in \eqref{eq:LB} forces $v = v_*$ and the operator
vanishes identically, so collisional relaxation appears only at higher order
in $1/N$. A one-dimensional run therefore tests the three-point cumulant
directly, but it is not a test of \eqref{eq:LB}.

\subsection{Retained state and cumulants}

\begin{itemize}
  \item $y_l$: $f_1$ in a weak basis, Fourier in $\bx$ and Hermite in $\bv$.
  \item $c_l$: the two-point cumulant in the same basis. Its lowest moment is
    the field fluctuation spectrum $\avg{|E_{\bk}|^2}$, which has a
    closed-form prediction from the dressed-particle picture at equilibrium.
    That prediction is the first sanity test.
  \item Scaling check: $\|g_2\|\sim N^{-1}$ and $\|g_3\|\sim N^{-2}$ across
    the sweep, with ensemble error bars.
\end{itemize}

\subsection{History object}

There is no collision log. The replacement is a per-layer momentum-transfer
record,
\begin{equation}
  \Delta\mathbf{p}_{ij}^{(l)} = -\frac{1}{N}\int_{t_l}^{t_{l+1}}
    \nabla\phi(\bx_i - \bx_j)\,\dd t ,
\end{equation}
kept as a sparse sketch, and its decomposition into three channels:
\begin{itemize}
  \item close encounters (impact parameter below the $90^\circ$ deflection
    radius), which can be logged as discrete events exactly as for hard
    spheres;
  \item weak binary deflections at intermediate range;
  \item collective transfer, measured in Fourier space at
    $k\lambda_D \lesssim 1$.
\end{itemize}
The fraction of momentum transfer in each channel, as a function of $\kappa$
and coupling, is the plasma counterpart of the molecule statistics, and it
locates a given system on the interpolation axis of the third target.

\subsection{Closure defects}

For a test function $\psi(\bx,\bv)$,
\begin{align}
  R^{\mathrm{V}}_\psi &= \frac{\dd}{\dd t}\avg{f_1,\psi}
     - \avg{f_1,\;\bv\cdot\nabla_{\bx}\psi
       - (\nabla\phi*\rho_f)\cdot\nabla_{\bv}\psi}, \\
  R^{\mathrm{LB}}_\psi &= R^{\mathrm{V}}_\psi
     - \frac{1}{N}\avg{Q_{\mathrm{LB}}(f_1,f_1),\psi} .
\end{align}
$R^{\mathrm{V}}$ is the Vlasov defect and should equal the collision term to
leading order; $R^{\mathrm{LB}}$ is the defect of the collisional closure and
should scale as $N^{-2}$ in the weakly coupled regime. Discretization error
and ensemble error must be reported separately from both.

\subsection{Arrow-of-time experiment}

The plasma version of the paired forward and reversed continuation is the
echo. At a layer boundary, reverse all velocities and continue. In the
Vlasov limit the phase-mixed perturbation reassembles; particle noise at
order $1/N$ degrades the echo. The echo amplitude as a function of $N$ and of
elapsed time measures directly how fast $g$ loses memory, which is the
assumption the Bogoliubov closure rests on. The three predictors of
\texttt{math\_to\_algorithm.md} (retained state; plus cumulant magnitudes;
plus time-oriented history) are compared on the echo data.

\subsection{Tests of the proposed mechanism}
\label{sec:tests}

The strategy of Sections~\ref{sec:scales}--\ref{sec:techniques} makes four
predictions that can be checked at times far shorter than the relaxation
time.
\begin{itemize}
  \item \textbf{Orbit decorrelation.} The departure of a tagged particle
    from its mean-field orbit satisfies
    $\avg{|\Delta x|^2} = \tfrac{2}{3}Dt^3$ with $D\propto 1/N$. This
    measures the Lenard--Balescu diffusion coefficient without a relaxation
    run.
  \item \textbf{Echo attenuation.} Echo amplitude as a function of delay $t$
    and of $N$ collapses onto a single curve in $k^2t^3/N$, with
    $\ln(\text{amplitude})$ linear in that variable.
  \item \textbf{History horizon.} History features older than
    $\tau_{\mathrm{mem}}\sim N^{1/3}$ add no predictive value for the
    defect; features inside the window may.
  \item \textbf{Channel fractions.} For a potential with short-range scale
    $b_{\min}$ and long-range scale $\lambda$, the fraction of velocity
    diffusion from non-overlapping binary encounters is
    $\ln(a/b_{\min})/\ln(\lambda/b_{\min})$, tending to $2/3$ for Coulomb in
    three dimensions. Finite-$N$ corrections are of relative order
    $1/\ln N$ and will be large in any affordable run.
\end{itemize}
Failure of the first two would refute the picture of
Section~\ref{sec:scales}; failure of the third would mean the memory cutoff
is not the relevant horizon for closure.

\paragraph{Preliminary check of the first prediction.} A scratch experiment,
not part of the repository and not a benchmark result, was run on
16 November 2026: particles on the two-dimensional torus $[0,2\pi)^2$ with
$\hat\phi(\bk) = |\bk|^{-2}$ truncated to $|k_i|\le 2$, uniform positions,
unit Maxwellian velocities, leapfrog with $\Delta t = 0.05$ to $t = 40$, four
realizations per $N$, fixed seed. The quantity $N\avg{|\Delta x|^2}$ for
$N = 500$, $2000$, $8000$ was
\begin{center}
\begin{tabular}{rrrrr}
  $t$ & $N=500$ & $N=2000$ & $N=8000$ & local slope ($N=8000$) \\ \hline
  1  & 1.57 & 1.55 & 1.95 & 3.50 \\
  5  & 260 & 237 & 260 & 2.98 \\
  10 & 1921 & 1789 & 2002 & 2.94 \\
  20 & 15236 & 14347 & 15766 & 3.00 \\
  40 & 119356 & 115393 & 127423 & 3.01 \\
\end{tabular}
\end{center}
The values agree across a factor of 16 in $N$ to within about 10\%, and the
logarithmic slope is 3 for $t\gtrsim 2$. No error bars were computed, so the
10\% spread is not separated from ensemble noise. This is consistent with
$\avg{|\Delta x|^2}\propto t^3/N$ and with a memory time
$\propto N^{1/3}$; it does not test the echo, the history horizon, or the
channel fractions.

\paragraph{Attempted check of the second prediction: no result.} An echo
experiment was run with the same system: velocity kicks
$0.05\cos x$ at $t=0$ and $0.05\cos 2x$ at $t=10$, looking for a response in
the $k=1$ density mode near $t=20$, with $N = 500$ to $32000$ and
$N\times\text{realizations} = 128000$. The measured peaks ($0.002$ to
$0.005$) were at the ensemble noise floor ($0.0028$), and a control with no
kicks gave the same peak ($0.005$). No echo was detected, so the run neither
supports nor refutes the predicted attenuation. A usable test needs a larger
echo (stronger kicks or weaker Landau damping of the pulses) or a much larger
ensemble.

\subsection{Limit-recovery tests}

\begin{itemize}
  \item Weak coupling, $N\to\infty$ at fixed time: Vlasov, with defect
    $\propto 1/N$.
  \item Weak coupling at time $N\tau$: relaxation rate against
    \eqref{eq:LB} and against Landau with a fitted Coulomb logarithm.
  \item Large $\kappa$, low density: crossover to Boltzmann-like sparse
    collision histories, where the existing hard-sphere track applies.
  \item Strong coupling: all three closures fail; the diagnostics should
    report it.
\end{itemize}

\section{Open questions}

\begin{enumerate}
  \item Is there a diagram expansion about the linearized Vlasov flow in
    which the non-ring diagrams can be counted and bounded in the way
    molecules are, with the gain coming from powers of $1/N$?
  \item What is the right layered ansatz? It must propagate cumulant size
    and enough velocity regularity to exclude echoes over the kinetic time.
  \item Does the three-channel decomposition of momentum transfer have a
    sharp formulation, so that the close-encounter channel can be treated by
    the hard-sphere method and the remainder by resummation?
  \item For the benchmark: do history features add predictive value for the
    Lenard--Balescu defect beyond $f_1$ and $g_2$, given that here $g_2$ is
    already part of the closure and not only its error?
\end{enumerate}

\begin{thebibliography}{99}
\small
\bibitem{DHM1} Y.~Deng, Z.~Hani, X.~Ma, \emph{Long time derivation of the
Boltzmann equation from hard sphere dynamics}, arXiv:2408.07818; to appear in
Ann.\ of Math.
\bibitem{DHM2} Y.~Deng, Z.~Hani, X.~Ma, \emph{Hilbert's sixth problem:
derivation of fluid equations via Boltzmann's kinetic theory},
arXiv:2503.01800 (preprint).
\bibitem{DH-WKE} Y.~Deng, Z.~Hani, \emph{Full derivation of the wave kinetic
equation}, Invent.\ Math.\ \textbf{233} (2023) 543--724.
\bibitem{DH-long} Y.~Deng, Z.~Hani, \emph{Long time justification of wave
turbulence theory}, arXiv:2311.10082 (preprint).
\bibitem{Lanford} O.~E. Lanford III, \emph{Time evolution of large classical
systems}, Lecture Notes in Physics \textbf{38}, Springer (1975).
\bibitem{GST} I.~Gallagher, L.~Saint-Raymond, B.~Texier, \emph{From Newton to
Boltzmann: Hard Spheres and Short-range Potentials}, Zurich Lectures in
Advanced Mathematics, EMS (2013); see also the published erratum.
\bibitem{BraunHepp} W.~Braun, K.~Hepp, \emph{The Vlasov dynamics and its
fluctuations in the $1/N$ limit of interacting classical particles}, Comm.\
Math.\ Phys.\ \textbf{56} (1977) 101--113.
\bibitem{Dobrushin} R.~L. Dobrushin, \emph{Vlasov equations}, Funct.\ Anal.\
Appl.\ \textbf{13} (1979) 115--123.
\bibitem{HaurayJabin} M.~Hauray, P.-E. Jabin, \emph{Particle approximation of
Vlasov equations with singular forces: propagation of chaos}, Ann.\ Sci.\
\'Ec.\ Norm.\ Sup\'er.\ (4) \textbf{48} (2015) 891--940.
\bibitem{LazaroviciPickl} D.~Lazarovici, P.~Pickl, \emph{A mean field limit
for the Vlasov--Poisson system}, Arch.\ Ration.\ Mech.\ Anal.\ \textbf{225}
(2017), no.~3.
\bibitem{Serfaty} S.~Serfaty, \emph{Mean field limit for Coulomb-type flows}
(appendix with M.~Duerinckx), Duke Math.\ J.\ \textbf{169} (2020) 2887--2935.
\bibitem{DuerinckxJabin} M.~Duerinckx, P.-E. Jabin, \emph{Derivation of 2D
Vlasov--Poisson for classical particles with Coulomb interactions},
arXiv:2608.04104 (2026, preprint, not refereed).
\bibitem{BPS} A.~V. Bobylev, M.~Pulvirenti, C.~Saffirio, \emph{From particle
systems to the Landau equation: a consistency result}, Comm.\ Math.\ Phys.\
\textbf{319} (2013) 683--702.
\bibitem{VelazquezWinter} J.~J.~L. Vel\'azquez, R.~Winter, \emph{From a
non-Markovian system to the Landau equation}, Comm.\ Math.\ Phys.\
\textbf{361} (2018) 239--287.
\bibitem{Winter} R.~Winter, \emph{Convergence to the Landau equation from the
truncated BBGKY hierarchy in the weak-coupling limit}, J.\ Differential
Equations \textbf{283} (2021) 1--36.
\bibitem{DuerinckxSR} M.~Duerinckx, L.~Saint-Raymond, \emph{Lenard--Balescu
correction to mean-field theory}, Probab.\ Math.\ Phys.\ \textbf{2} (2021)
27--69. The admissible exponent $r$ is stated differently in the abstract and
in the arXiv theorem; check the published Theorem~1 before quoting a value.
\bibitem{DuerinckxChaos} M.~Duerinckx, \emph{On the size of chaos via Glauber
calculus in the classical mean-field dynamics}, Comm.\ Math.\ Phys.\
\textbf{382} (2021) 613--653.
\bibitem{DuerinckxLeBihan} M.~Duerinckx, C.~Le~Bihan, \emph{Lenard--Balescu
thermalization: rigorous derivation from a toy model}, arXiv:2511.10778
(2025, preprint).
\bibitem{MouhotVillani} C.~Mouhot, C.~Villani, \emph{On Landau damping}, Acta
Math.\ \textbf{207} (2011) 29--201.
\bibitem{BMM} J.~Bedrossian, N.~Masmoudi, C.~Mouhot, \emph{Landau damping:
paraproducts and Gevrey regularity}, Ann.\ PDE \textbf{2} (2016), Art.~4.
\bibitem{BedrossianEcho} J.~Bedrossian, \emph{Nonlinear echoes and Landau
damping with insufficient regularity}, Tunis.\ J.\ Math.\ \textbf{3} (2021)
121--205.
\bibitem{GuoJang} Y.~Guo, J.~Jang, \emph{Global Hilbert expansion for the
Vlasov--Poisson--Boltzmann system}, Comm.\ Math.\ Phys.\ \textbf{299} (2010)
469--501.
\bibitem{ArsenioSR} D.~Ars\'enio, L.~Saint-Raymond, \emph{From the
Vlasov--Maxwell--Boltzmann System to Incompressible Viscous
Electro-magneto-hydrodynamics, Volume~1}, EMS Monographs in Mathematics
(2019).
\bibitem{BardosBesse} C.~Bardos, N.~Besse, \emph{Diffusion limit of the
Vlasov equation in the weak turbulent regime}, J.\ Math.\ Phys.\ \textbf{62}
(2021) 101505.
\bibitem{Dawson} J.~Dawson, \emph{Thermal relaxation in a one-species,
one-dimensional plasma}, Phys.\ Fluids \textbf{7} (1964) 419--425.
\bibitem{Fouvry} J.-B. Fouvry, B.~Bar-Or, P.-H. Chavanis, \emph{Kinetic theory
of one-dimensional homogeneous long-range interacting systems sourced by
$1/N^2$ effects}, Phys.\ Rev.\ E \textbf{100} (2019) 052142.
\bibitem{BodineauLeBris} T.~Bodineau, P.~Le~Bris, \emph{Linear Landau
equation as a limit of a tagged particle in mean field interaction with a
free gas}, arXiv:2602.16440 (2026, preprint).
\bibitem{Bedrossian17} J.~Bedrossian, \emph{Suppression of plasma echoes and
Landau damping in Sobolev spaces by weak collisions in a
Vlasov--Fokker--Planck equation}, Ann.\ PDE \textbf{3} (2017), Art.~19.
\bibitem{CLN} S.~Chaturvedi, J.~Luk, T.~T. Nguyen, \emph{The
Vlasov--Poisson--Landau system in the weakly collisional regime}, J.\ Amer.\
Math.\ Soc.\ \textbf{36} (2023) 1103--1189.
\bibitem{ZhaoZi} W.~Zhao, R.~Zi, \emph{Optimal stability threshold in lower
regularity spaces for the Vlasov--Poisson--Fokker--Planck equations},
arXiv:2603.29204 (2026, preprint; abstract only consulted).
\bibitem{Dupree} T.~H. Dupree, \emph{A perturbation theory for strong plasma
turbulence}, Phys.\ Fluids \textbf{9} (1966) 1773--1782 (consulted through
secondary sources).
\bibitem{ONeil} T.~M. O'Neil, \emph{Effect of Coulomb collisions and
microturbulence on the plasma wave echo}, Phys.\ Fluids \textbf{11} (1968)
2420--2425 (consulted through secondary sources).
\bibitem{NVW} A.~Nota, J.~J.~L. Vel\'azquez, R.~Winter, \emph{Interacting
particle systems with long-range interactions: scaling limits and kinetic
equations}, Atti Accad.\ Naz.\ Lincei Rend.\ Lincei Mat.\ Appl.\ \textbf{32}
(2021) 335--377.
\bibitem{Lancellotti} C.~Lancellotti, \emph{On the fluctuations about the
Vlasov limit for $N$-particle systems with mean-field interactions}, J.\
Stat.\ Phys.\ \textbf{136} (2009) 643--665.
\bibitem{FeliachiBouchet} O.~Feliachi, F.~Bouchet, \emph{Dynamical large
deviations for homogeneous systems with long range interactions and the
Balescu--Guernsey--Lenard equation}, arXiv:2105.05644.
\bibitem{CagliotiRousset} E.~Caglioti, F.~Rousset, \emph{Quasi-stationary
states for particle systems in the mean-field limit}, J.\ Stat.\ Phys.\
\textbf{129} (2007) 241--263.
\end{thebibliography}

\end{document}
